% Chemistry Homework Formatter — student paper preamble
% Default wrapper is English (CIE papers). Load ctex only if the teacher asks for Chinese.
% Compiler: XeLaTeX. Overleaf: upload this file as main.tex plus the images/ folder.
\documentclass[10pt,a4paper]{article}

\usepackage[margin=16mm,headheight=13pt,headsep=5pt,footskip=9mm]{geometry}
\usepackage{fontspec}
\usepackage[version=4]{mhchem}
\usepackage{chemfig}
\usepackage{graphicx}
\usepackage{xcolor}
\usepackage{booktabs}
\usepackage{array}
\usepackage{enumitem}
\usepackage{needspace}
\usepackage{setspace}
\usepackage{fancyhdr}
\usepackage{amsmath,amssymb}

\raggedbottom
\setlist[enumerate]{leftmargin=1.6em,itemsep=0.25em,topsep=0.25em}
\setlist[itemize]{leftmargin=1.4em,itemsep=0.12em}

\pagestyle{fancy}
\fancyhf{}
\fancyhead[L]{\small \homeworktitle}
\fancyhead[R]{\small \homeworkmarks}
\fancyfoot[C]{\thepage}
\renewcommand{\headrulewidth}{0.3pt}

\newcommand{\homeworktitle}{Chemistry homework}
\newcommand{\homeworkmarks}{}
\newcommand{\sethomework}[2]{%
  \renewcommand{\homeworktitle}{#1}%
  \renewcommand{\homeworkmarks}{#2}%
}

\newcommand{\answerlines}[1]{%
  \par\vspace{0.3em}%
  \foreach \n in {1,...,#1}{\noindent\rule{\linewidth}{0.2pt}\\[1.05em]}%
}

% Fallback if pgffor is missing: comment \answerlines and use \vspace instead.
\usepackage{pgffor}

\graphicspath{{images/}}
\DeclareGraphicsExtensions{.png,.jpg,.jpeg,.pdf}

% Measure the image, then break only if the whole labelled clip will not fit.
\newcommand{\clipq}[3]{%
  \par
  \sbox0{\includegraphics[width=0.94\linewidth,keepaspectratio]{#3}}%
  \Needspace{\dimexpr\ht0+1.35em+2pt\relax}%
  \noindent{\small\textbf{#1}}\hfill{\small[#2]}\par\vspace{0.08em}%
  \noindent\usebox0\par
  \vspace{0.28em}%
}

\newcommand{\qfigure}[1]{%
  \sbox0{\includegraphics[width=0.94\linewidth,keepaspectratio]{#1}}%
  \Needspace{\dimexpr\ht0+0.6em\relax}%
  \noindent\usebox0\par
}

\newcommand{\qfigurebig}[1]{%
  \sbox0{\includegraphics[width=\linewidth,keepaspectratio]{#1}}%
  \Needspace{\dimexpr\ht0+0.6em\relax}%
  \noindent\usebox0\par
}

\setlength{\parindent}{0pt}
\setlength{\parskip}{0.12em}

\setlist[enumerate]{leftmargin=1.55em,itemsep=0.06em,topsep=0.12em}
\newcommand{\qhead}[1]{%
  \par\vspace{0.38em}%
  {\normalsize\bfseries #1}\par\vspace{0.12em}%
}
\newcommand{\mcq}[1]{%
  \par\noindent\textbf{#1.}\enspace\ignorespaces
}
\newcommand{\mcqend}{\hfill{[1]}\par}
\newcommand{\optrow}[4]{%
\par\noindent{\small
\begin{tabular}{@{}p{0.48\linewidth}@{}p{0.48\linewidth}@{}}
\textbf{A}\; #1 & \textbf{B}\; #2 \\[0.12em]
\textbf{C}\; #3 & \textbf{D}\; #4
\end{tabular}}\par
}
\newcommand{\optlist}[4]{%
\par
\begin{enumerate}[label=\textbf{\Alph*.}, leftmargin=1.55em, itemsep=0.05em, topsep=0.08em]
\item #1
\item #2
\item #3
\item #4
\end{enumerate}
}
\newcommand{\qfig}[2][0.58]{%
\par
\begin{center}\vspace{-0.1em}%
\includegraphics[width=#1\linewidth,keepaspectratio]{#2}%
\end{center}\vspace{-0.1em}%
}

\begin{document}
\sethomework{AS stoichiometry homework 1 · mole foundations}{14 marks}

\begin{center}
{\large\bfseries AS stoichiometry homework 1 · mole foundations}\\[0.2em]
{\small CIE 9701 AS Chemistry · 14 MCQ · 14 marks}
\end{center}

\noindent Name: \underline{\hspace{5.8cm}} \quad Class: \underline{\hspace{2.6cm}} \quad Date: \underline{\hspace{2.8cm}}

\vspace{0.25em}
{\small Topics: 1 masses of atoms and molecules · 2 the mole and Avogadro · 3 reacting masses · 4 complete combustion.

Circle \textbf{one} letter (A--D) for each question. Show working. Calculators and a Periodic Table may be used.}
\vspace{0.15em}

\noindent\begin{minipage}{\linewidth}
\qhead{1 · Masses of atoms and molecules}
\mcq{1}%
The mass spectrum of a sample of neon is shown. The relative abundance of each peak is written in brackets above it.
\qfig[0.48]{fig-q01-spectrum.png}
What is the relative atomic mass, $A_\mathrm{r}$, of this sample of neon?
\optrow{20.15}{20.20}{21.00}{21.82}
\mcqend
\end{minipage}\vspace{0.16em}

\noindent\begin{minipage}{\linewidth}
\mcq{2}%
A sample of sulfur consists mostly of \ce{^{32}S}. It also contains 4.2\% \ce{^{34}S} and 2.8\% \ce{^{36}S}. No other isotopes of sulfur are present.

What is the relative atomic mass, $A_\mathrm{r}$, of \textbf{this} sample of sulfur?
\optrow{32.1}{32.2}{34.0}{34.3}
\mcqend
\end{minipage}\vspace{0.16em}

\noindent\begin{minipage}{\linewidth}
\mcq{3}%
The diagram shows the apparatus used to find the relative molecular mass of a volatile liquid.

When 0.10\,g of a volatile liquid is injected into the syringe, all of the volatile liquid evaporates and the volume increases by 85\,cm$^3$.

The heater maintains a temperature of 400\,K and the experiment is carried out at a pressure of 101\,300\,Pa.
\qfig[0.78]{fig-q03-apparatus.png}
If the vapour of the volatile liquid behaves as an ideal gas, which expression can be used to calculate the relative molecular mass of the liquid?
\optlist{$M_\mathrm{r}=(85\times 101\,300)\div(0.10\times 8.31\times 400)$}{$M_\mathrm{r}=(85\times 101.3)\div(0.10\times 8.31\times 400)$}{$M_\mathrm{r}=(0.10\times 8.31\times 400)\div(85\times 10^{-6}\times 101\,300)$}{$M_\mathrm{r}=(0.10\times 8.31\times 400)\div(85\times 10^{-6}\times 101.3)$}
\mcqend
\end{minipage}\vspace{0.16em}

\noindent\begin{minipage}{\linewidth}
\qhead{2 · The mole and the Avogadro constant}
\mcq{4}%
Which statement about the Avogadro constant is correct?
\optlist{It is the mass of one mole of any element.}{It is the mass of $6.02\times 10^{23}$ atoms of any element.}{It is the number of atoms in one mole of neon.}{It is the number of atoms in 12\,g of any element.}
\mcqend
\end{minipage}\vspace{0.16em}

\noindent\begin{minipage}{\linewidth}
\mcq{5}%
Which contains the largest number of hydrogen atoms?
\optlist{0.10\,mol of pentane}{0.20\,mol of but-2-ene}{1.00\,mol of hydrogen molecules}{$6.02\times 10^{23}$ hydrogen atoms}
\mcqend
\end{minipage}\vspace{0.16em}

\noindent\begin{minipage}{\linewidth}
\mcq{6}%
Which sample contains the most iodine?
\optrow{1\,g of \ce{CaI2}}{1\,g of \ce{KI}}{1\,g of \ce{NaI}}{1\,g of \ce{NH4I}}
\mcqend
\end{minipage}\vspace{0.16em}

\noindent\begin{minipage}{\linewidth}
\mcq{7}%
What contains $9.03\times 10^{23}$ oxygen atoms?
\optlist{0.25\,mol aluminium oxide}{0.75\,mol sulfur dioxide}{1.5\,mol sulfur trioxide}{3.0\,mol water}
\mcqend
\end{minipage}\vspace{0.16em}

\noindent\begin{minipage}{\linewidth}
\qhead{3 · Reacting masses}
\mcq{8}%
Sodium peroxide, \ce{Na2O2}, is used to absorb carbon dioxide from the atmosphere and release oxygen in closed environments such as space capsules and submarines.
\[\ce{2Na2O2 + 2CO2 -> 2Na2CO3 + O2}\]
Which mass of sodium peroxide would be required to remove 2.4\,dm$^3$ of carbon dioxide from the atmosphere at room temperature and pressure?
\optrow{2.4\,g}{3.9\,g}{7.8\,g}{15.6\,g}
\mcqend
\end{minipage}\vspace{0.16em}

\noindent\begin{minipage}{\linewidth}
\mcq{9}%
A student reacts 0.100\,mol of each of sodium, magnesium and phosphorus atoms separately with an excess of oxygen.

Which rows are correct?

{\centering
{\small
\begin{tabular}{c l c}
\toprule
 & oxide & mass of oxide formed\,/\,g \\
\midrule
1 & sodium & 3.10 \\
2 & magnesium & 4.03 \\
3 & phosphorus & 7.10 \\
\bottomrule
\end{tabular}}\par}

\optrow{1, 2 and 3}{1 and 2 only}{1 and 3 only}{2 and 3 only}
\mcqend
\end{minipage}\vspace{0.16em}

\noindent\begin{minipage}{\linewidth}
\mcq{10}%
Separate samples, each of mass 1.0\,g, of the compounds listed are treated with an excess of dilute acid.

Which compound releases the largest amount of \ce{CO2}?
\optrow{1.0\,g \ce{CaCO3}}{1.0\,g \ce{Li2CO3}}{1.0\,g \ce{MgCO3}}{1.0\,g \ce{Na2CO3}}
\mcqend
\end{minipage}\vspace{0.16em}

\noindent\begin{minipage}{\linewidth}
\mcq{11}%
2.0\,g of ammonium nitrate, \ce{NH4NO3}, decomposes to give 0.90\,g of water and a single gas.

What is the identity of the gas?
\optrow{\ce{NO}}{\ce{NO2}}{\ce{N2O}}{\ce{N2}}
\mcqend
\end{minipage}\vspace{0.16em}

\noindent\begin{minipage}{\linewidth}
\qhead{4 · Complete combustion}
\mcq{12}%
What is the minimum mass of oxygen required to ensure the complete combustion of 12\,dm$^3$ of propane measured under room conditions?
\optrow{60\,g}{80\,g}{120\,g}{160\,g}
\mcqend
\end{minipage}\vspace{0.16em}

\noindent\begin{minipage}{\linewidth}
\mcq{13}%
Mixture R consists of one mole of \ce{C3H6} and one mole of \ce{C4H6}.

What is the minimum number of moles of oxygen molecules needed for complete combustion of mixture R?
\optrow{6.5}{7}{10}{20}
\mcqend
\end{minipage}\vspace{0.16em}

\noindent\begin{minipage}{\linewidth}
\mcq{14}%
If 1 mole of hexane combusts in an excess of oxygen, how many moles of products are formed?
\optrow{11}{12}{13}{14}
\mcqend
\end{minipage}\vspace{0.16em}

\end{document}
